\chapter{The Laplace--Beltrami Operator and Mesh Smoothing}\label{ch:smooth}

Scanned meshes are noisy: measurement errors move every vertex a little along and across the
surface. Smoothing removes this noise while keeping the shape. The tool for the job is the
Laplace--Beltrami operator, which plays on surfaces the role that the second derivative plays on
the real line. This chapter derives its discretisations, studies its eigenfunctions as a Fourier
basis, and analyses the diffusion flows that smooth a mesh. It follows the slides of Tema 4 and the
companion notebook. We use $\Lap=\Div\grad$ (negative semi-definite) and its discretisation
$\mM^{-1}\mC$; the eigenvalues are written as positive numbers $\lambda$ with
$-\mC\bphi=\lambda\mM\bphi$.

\section{Smoothing as filtering}

On the circle, a function $f(x)=\sum_kc_ke^{\mathrm ikx}$ is a sum of oscillations of frequency
$k$, and each one is an eigenfunction of the second derivative:
$\Lap e^{\mathrm ikx}=-k^2e^{\mathrm ikx}$. The Laplacian therefore multiplies high frequencies by
large factors. The \emph{heat equation} $\partial_tf=\Lap f$ has the solution
\begin{equation}
f(x,t)=\sum_kc_ke^{-k^2t}e^{\mathrm ikx}:
\end{equation}
every mode decays at a rate equal to its eigenvalue, so diffusion is a low-pass filter whose cut-off
frequency decreases with time (\cref{fig:heat1d}). The noise of a scan lives in high frequencies
and the shape in low frequencies, so we need the analogue of $\Lap$, of its eigenfunctions and of
the heat equation on a surface.

\fig[0.85\textwidth]{ch05_heat1d.pdf}{Left: for $f=\sin3x+0.3\sin25x$, the Laplacian multiplies
the slow mode by $9$ and the fast one by $625$, so $\Lap f$ (scaled by $1/400$) is dominated by
the fast mode. Right: the heat equation damps every mode as $e^{-k^2t}$, removing the fast
oscillation long before the slow one changes.}{fig:heat1d}

\section{The Laplace--Beltrami operator}

On a surface $S$, the gradient $\grad_Sf$ of a function is the tangent vector field that gives the
directional derivatives, $\dd f(\bw)=\inner{\grad_Sf}{\bw}$; the divergence of a tangent field
measures its net outflow per unit area; and the Laplace--Beltrami operator is
$\Lap_Sf=\Div_S\grad_Sf$. In a parameterization with metric $\mathbf I=(g_{ij})$,
\begin{equation}
\Lap_Sf=\frac{1}{\sqrt{\det g}}\sum_{i,j}\partial_i\!\left(\sqrt{\det g}\,g^{ij}\partial_jf\right),
\end{equation}
which depends only on the metric, so $\Lap_S$ is intrinsic. Its key properties are:
\begin{itemize}
  \item $\Lap_S$ annihilates constants, and on a closed surface only constants are harmonic;
  \item Green's identity $\int_Sg\Lap_Sf\dd A=-\int_S\inner{\grad_Sf}{\grad_Sg}\dd A$ makes it
  symmetric and negative semi-definite;
  \item it is local: $(\Lap_Sf)(\bx)$ depends only on $f$ near $\bx$;
  \item applied to the coordinates it gives the mean curvature normal, $\Lap_S\bx=-2H\bn$
  (\cref{ch:diffgeo}).
\end{itemize}

\section{Discretisations}

\subsection{The uniform Laplacian}

The simplest discretisation uses only the connectivity:
\begin{equation}
(\mL_u\bfun)_i=\frac1{d_i}\sum_{j\in\mathcal N(i)}(f_j-f_i),
\end{equation}
the difference between the average of the neighbours and the vertex itself. In matrix form
$\mL_u=\mD^{-1}\mA-\mI$, with $\mA$ the adjacency matrix and $\mD$ the diagonal matrix of
valences. Applied to positions, $\mL_u\bx$ points from the vertex to the centroid of its
one-ring. On a regular grid it is the classical five-point stencil, but on irregular meshes the
centroid differs from the vertex even when the surface is flat, so $\mL_u\bx$ has a
\emph{tangential} component that moves vertices along the surface (\cref{fig:flat}), and $\mL_u$
does not converge to $\Lap_S$ under refinement.

\subsection{The cotangent Laplacian by finite elements}

Let $\varphi_i$ be the piecewise linear ``hat'' function equal to $1$ at vertex $i$ and $0$ at
every other vertex, and represent $f=\sum_jf_j\varphi_j$. The Laplacian of a piecewise linear
function is zero inside each triangle and singular on the edges, so we discretise the weak form
instead: multiply $\Lap f=g$ by a test function $\varphi_i$, integrate and apply Green's identity,
\begin{equation}
\int_S\varphi_i\,g\dd A=-\int_S\inner{\grad\varphi_i}{\grad f}\dd A .
\end{equation}
Expanding $f$ and $g$ in the hat basis gives the linear system
\begin{equation}
\mM\bg=\mC\bfun,\qquad M_{ij}=\int_S\varphi_i\varphi_j\dd A,\qquad C_{ij}=-\int_S\inner{\grad\varphi_i}{\grad\varphi_j}\dd A .
\end{equation}

\paragraph{The stiffness matrix.} In a triangle with corners $i,j,k$ and area $A$, the gradient of
$\varphi_i$ is constant, orthogonal to the opposite edge $\be_{jk}=\bx_k-\bx_j$, and of length
$1/h_i$, where $h_i$ is the height from $i$: $\grad\varphi_i=(\bn\times\be_{jk})/(2A)$. Then
\begin{equation}
\int_T\inner{\grad\varphi_i}{\grad\varphi_j}\dd A=\frac{\inner{\be_{jk}}{\be_{ki}}}{4A}
=-\frac{\norm{\be_{jk}}\norm{\be_{ki}}\cos\theta_k}{2\norm{\be_{jk}}\norm{\be_{ki}}\sin\theta_k}=-\frac12\cot\theta_k ,
\end{equation}
using $2A=\norm{\be_{jk}}\norm{\be_{ki}}\sin\theta_k$, where $\theta_k$ is the angle at the third
corner. Summing over the two triangles that share the edge $ij$, with opposite angles
$\alpha_{ij}$ and $\beta_{ij}$:
\begin{equation}
C_{ij}=\tfrac12(\cot\alpha_{ij}+\cot\beta_{ij})\quad(j\in\mathcal N(i)),\qquad C_{ii}=-\sum_{j\neq i}C_{ij},
\label{eq:cotan-matrix}
\end{equation}
and $C_{ij}=0$ otherwise. $\mC$ is symmetric, sparse (about seven non-zeros per row), its rows sum to
zero, and it is negative semi-definite.

\paragraph{The mass matrix.} The exact (consistent) mass matrix has entries $A/6$ on the diagonal
and $A/12$ off the diagonal per triangle. In practice it is \emph{lumped} into a diagonal matrix
$\mM=\diag(A_1,\dots,A_n)$ with the barycentric or mixed Voronoi vertex areas of \cref{ch:diffgeo}.
The discrete Laplace--Beltrami operator is then
\begin{equation}
\Lap\approx\mM^{-1}\mC,\qquad (\Lap\bfun)_i=\frac{1}{2A_i}\sum_{j\in\mathcal N(i)}(\cot\alpha_{ij}+\cot\beta_{ij})(f_j-f_i),
\end{equation}
which is the cotangent formula \eqref{eq:cotan} of \cref{ch:diffgeo}. Note that
$\mM^{-1}\mC$ is not symmetric; algorithms work with the symmetric pair $(\mC,\mM)$ instead.

\begin{example}{Cotangent weights}{cotw}
Consider the edge $ij$ shared by two triangles whose angles opposite to it are $\alpha=60^\circ$
and $\beta=90^\circ$. Then $C_{ij}=\tfrac12(\cot60^\circ+\cot90^\circ)=\tfrac12(0.577+0)=0.289$.
If instead $\beta=120^\circ$, $\cot\beta=-0.577$ and $C_{ij}=0$: the two contributions cancel.
For $\alpha+\beta>\pi$ the weight is negative. Since
$\cot\alpha+\cot\beta=\sin(\alpha+\beta)/(\sin\alpha\sin\beta)$, the weight is non-negative exactly
when $\alpha+\beta\le\pi$, the local \emph{Delaunay} condition for the edge.
\end{example}

\subsection{Comparing the discretisations}

The cotangent Laplacian has \emph{linear precision}: if the mesh is planar and $f$ is a linear
function, then $\mC\bfun=\bzero$ at interior vertices, exactly as $\Lap f=0$ in the continuous
setting. Applied to the coordinates of a flat mesh, $\mC\bx=\bzero$, so cotangent smoothing never
moves vertices within a plane, whereas the uniform Laplacian does (\cref{fig:flat}). It also
converges to $\Lap_S$ under refinement of reasonable meshes and depends only on the geometry, not
on how the mesh was triangulated. Its drawback is that weights become negative on non-Delaunay
edges, which can break the maximum principle.

\fig[0.7\textwidth]{ch05_flat_uniform.pdf}{An irregular triangulation of a planar square. Left:
the uniform Laplacian of the vertex positions (red arrows) is far from zero; its largest interior
norm is $0.0515$, so uniform smoothing slides vertices inside the plane. Right: the cotangent
Laplacian vanishes at every interior vertex up to round-off ($3.6\times10^{-15}$).}{fig:flat}

\Citet{Wardetzky2007} proved that on general meshes no discrete Laplacian can be symmetric,
local, have linear precision and have non-negative weights at the same time (``no free lunch''). The
uniform Laplacian sacrifices linear precision, the cotangent Laplacian sacrifices positivity on
non-Delaunay meshes. Intrinsic Delaunay triangulations, and the robust Laplacian of
\citet{Sharp2020} for non-manifold meshes and point clouds, recover positivity without moving
vertices.

\section{The spectral view}

The discrete eigenproblem
\begin{equation}
-\mC\bphi_k=\lambda_k\mM\bphi_k,\qquad 0=\lambda_1<\lambda_2\le\lambda_3\le\cdots,\qquad
\bphi_k\T\mM\bphi_l=\delta_{kl},
\label{eq:geneig}
\end{equation}
is a generalized symmetric eigenproblem whose eigenvectors are orthonormal with respect to the
$\mM$ inner product $\inner{\bfun}{\bg}_\mM=\bfun\T\mM\bg\approx\int_Sfg\dd A$. The first
eigenvector is constant. On the unit sphere the eigenfunctions are the spherical harmonics, with
eigenvalue $\ell(\ell+1)$ of multiplicity $2\ell+1$; the discrete spectrum reproduces these
plateaus (\cref{fig:sphere-spec}). On a general shape the low eigenfunctions vary slowly along the
main parts of the object and the higher ones oscillate with increasing frequency
(\cref{fig:eigvec}); they are the \emph{manifold harmonics} \citep{Vallet2008,Zhang2010}.

\fign{ch05_sphere_spectrum.pdf}{First eigenvalues of the cotangent Laplacian on an irregularly
triangulated unit sphere. They group in plateaus of $2\ell+1$ values at
$\ell(\ell+1)=0,2,6,12,20,\dots$ (dotted lines), the eigenvalues of the spherical
harmonics.}{fig:sphere-spec}

\fig{ch05_eigenvectors.png}{Eigenvectors of the cotangent Laplacian on the bunny (blue negative,
red positive). Low frequencies follow the large-scale structure (the ears, the body); the index
and the eigenvalue grow with the number of oscillations.}{fig:eigvec}

Any function, and in particular each coordinate of the vertex positions, can be expanded in this
basis. The coefficients are $\bc=\mPhi\T\mM\bx$ (one row per eigenvector, one column per
coordinate) and the \emph{spectral low-pass filter} keeps the first $m$ terms,
\begin{equation}
\tilde\bx=\sum_{k=1}^{m}\bphi_k\bphi_k\T\mM\bx=\mPhi_m\mPhi_m\T\mM\bx .
\end{equation}
\Cref{fig:lowpass} shows reconstructions of a noisy bunny. The coefficients of the clean shape
decay quickly with $\lambda$ while the noise spreads evenly over all frequencies, so truncation
separates them. Spectral filtering is the cleanest form of smoothing, but computing many
eigenvectors is expensive: the first few hundred are obtained with shift-invert Lanczos iterations
(\texttt{scipy.sparse.linalg.eigsh} with a small negative shift), each requiring a sparse
factorisation; the full spectrum needs dense $O(n^3)$ methods and is only practical for small
meshes. Diffusion flows achieve a similar filtering with a few sparse solves.

\fig{ch05_lowpass.png}{Spectral low-pass filtering of a noisy bunny with the first $m$
eigenvectors. With $m=10$ only the coarse shape survives, $m=150$ already shows the ears and
legs, and $m=600$ starts to bring back noise. Right: norm of the spectral coefficients of the noisy
and clean shapes against $\lambda_k$: the signal decays, the noise is flat.}{fig:lowpass}

\section{Diffusion flow}

Smoothing a mesh by diffusion means evolving the vertex positions with the heat equation
\begin{equation}
\frac{\partial\bx}{\partial t}=\Lap\bx ,
\end{equation}
which, by \eqref{eq:mean-curvature-normal}, moves each point against the normal with speed $2H$.
In the eigenbasis the exact solution for a fixed operator is $\bc_k(t)=e^{-\lambda_kt}\bc_k(0)$:
diffusion for time $t$ is a low-pass filter with transfer function $e^{-t\lambda}$. Two time
discretisations are used.

\subsection{Explicit Euler}

The explicit (forward Euler) step with step size $h$ is
\begin{equation}
\bx^{(k+1)}=\bx^{(k)}+h\,\Lap\bx^{(k)} .
\end{equation}
With the uniform Laplacian and $h=\lambda\in(0,1)$ this is classical \emph{Laplacian smoothing}:
every vertex moves a fraction $\lambda$ of the way towards the centroid of its neighbours. In the
eigenbasis each step multiplies mode $k$ by $(1-h\lambda_k)$, so the iteration is stable only if
$\abs{1-h\lambda_k}\le1$ for every $k$:
\begin{equation}
h\le\frac{2}{\lambda_{\max}} .
\end{equation}
For the uniform Laplacian $\lambda_{\max}\le2$, so $h<1$. For the cotangent Laplacian $\lambda_{\max}$
grows like $1/\bar\ell^2$, where $\bar\ell$ is the edge length, and a single small or sliver triangle
can make it huge, so explicit cotangent smoothing needs tiny steps. Close to the limit the highest
modes are not damped but flip sign at every step, and beyond it they grow exponentially
(\cref{fig:transfer}).

\subsection{Implicit Euler}

The implicit (backward Euler) step evaluates the Laplacian at the new positions,
$\bx^{(k+1)}=\bx^{(k)}+h\,\mM^{-1}\mC\bx^{(k+1)}$, that is
\begin{equation}
(\mM-h\mC)\,\bx^{(k+1)}=\mM\bx^{(k)} .
\label{eq:implicit}
\end{equation}
The matrix $\mM-h\mC$ is sparse, symmetric and positive definite, so the system is solved with a
sparse Cholesky factorisation (reused for all three coordinates and for all steps if the operator is
kept fixed) or with conjugate gradients. Each step multiplies mode $k$ by $1/(1+h\lambda_k)$, which
is in $(0,1]$ for every $h>0$: the scheme is \emph{unconditionally stable} \citep{Desbrun1999}. A
single implicit step of size $t$ already approximates the ideal filter $e^{-t\lambda}$, and $N$
steps of size $t/N$ approach it as $(1+t\lambda/N)^{-N}\to e^{-t\lambda}$
(\cref{fig:transfer}).

\fig{ch05_transfer.pdf}{Transfer functions of smoothing schemes. Left: for total time $t$, the
exact filter $e^{-t\lambda}$, one implicit step $1/(1+t\lambda)$, ten implicit steps, and an
explicit scheme near its stability limit, whose transfer oscillates in sign for large $\lambda$.
Centre: after many explicit steps, $h=1.9/\lambda_{\max}$ still damps all modes, while
$h=2.1/\lambda_{\max}$ amplifies the highest ones without bound. Right: ten steps of Laplacian
smoothing with $\lambda=0.33$ against Taubin's $\lambda|\mu$ filter, which passes the lowest
frequencies almost unchanged.}{fig:transfer}

\paragraph{Linear versus curvature flow.} If $\mC$ and $\mM$ are computed once on the input mesh,
the flow is linear and every step reuses the same factorisation. Recomputing them at every step
gives a discretisation of the true \emph{mean curvature flow}, which is non-linear and more
faithful to the geometry but more expensive.

\section{Shrinkage, drift and their remedies}

Diffusion damps every non-constant mode, including the low frequencies that carry the shape. In
the limit $t\to\infty$ every coordinate converges to its mean and the mesh collapses to a point.
Under mean curvature flow a sphere of radius $r_0$ shrinks as $r(t)=\sqrt{r_0^2-4t}$ and
disappears at $t=r_0^2/4$. Three remedies are standard.

\paragraph{Volume rescaling.} The enclosed volume of a closed oriented mesh follows from the
divergence theorem,
\begin{equation}
\vol=\frac16\sum_{f=(a,b,c)}\bx_a\cdot(\bx_b\times\bx_c),
\end{equation}
and after each step the mesh is scaled about its centroid by $(\vol_0/\vol)^{1/3}$ to restore the
initial volume \citep{Desbrun1999}.

\paragraph{Taubin's $\lambda|\mu$ smoothing.} \citet{Taubin1995} alternates a shrinking step with
$\lambda>0$ and an inflating step with $\mu<-\lambda<0$:
\begin{equation}
\bx\leftarrow(\mI+\lambda\mL_u)\bx,\qquad\bx\leftarrow(\mI+\mu\mL_u)\bx .
\end{equation}
For an eigenvalue $\kappa$ of $-\mL_u$ the transfer of one pair of steps is
$f(\kappa)=(1-\lambda\kappa)(1-\mu\kappa)$. It equals one at $\kappa=0$ and at the
\emph{pass-band frequency} $\kappa_{PB}=1/\lambda+1/\mu>0$, is slightly larger than one in between
and smaller than one above. With the usual values $\lambda=0.33$, $\mu=-0.34$, $\kappa_{PB}\approx0.09$:
low frequencies are preserved and high frequencies removed, so the filter smooths without shrinking
(\cref{fig:transfer}, right).

\paragraph{Using the right operator.} The uniform Laplacian also has a tangential component that
drags vertices towards densely sampled regions. \Cref{fig:irregular} smooths a sphere whose
triangulation is denser on one side until its area halves: with the uniform operator
the centre moves by $0.232$ and the radius varies by $5.7\%$ over the surface, while the cotangent
operator keeps the centre and the spherical shape ($0.1\%$ deviation).

\fig[0.7\textwidth]{ch05_irregular_sphere.png}{Smoothing an irregularly sampled unit sphere until
its area halves. The uniform Laplacian shifts the centroid by $0.232$ and deforms the sphere
($5.7\%$ radial deviation); the cotangent Laplacian shrinks it uniformly ($0.1\%$
deviation).}{fig:irregular}

\section{Putting it together}

\Cref{fig:implicit} shows implicit cotangent smoothing of a noisy bunny with growing step sizes,
measured in units of $\bar\ell^2$ (the squared mean edge length, the natural time scale of
diffusion over one edge). Small steps remove the noise; large steps remove the ears, the legs and
finally the bunny itself, which degenerates into a thin, elongated blob before collapsing. In
practice smoothing must be stopped: \cref{fig:errcurves} measures the distance to the clean
surface as smoothing proceeds. The error first drops below the noise level, reaches a minimum and
then grows again as shrinkage dominates; volume rescaling lowers the minimum and widens the
useful range of $t$.

\fig{ch05_implicit.png}{Implicit cotangent smoothing \eqref{eq:implicit} of a noisy bunny with one
step of size $h=1,10,100,1000\,\bar\ell^2$. The last panel shows the result for $h=1000\bar\ell^2$
rescaled to the original size: the shape has collapsed towards a curve.}{fig:implicit}

\fig[0.85\textwidth]{ch05_error_curves.pdf}{Median distance to the clean surface while smoothing
a noisy test surface (dotted line: noise level). Left: explicit uniform smoothing with $h=0.5$
reaches its best result after a few iterations and then shrinks the shape. Right: implicit
cotangent smoothing with small steps, the same with volume rescaling after every step, and single
implicit steps of size $t$.}{fig:errcurves}

A practical recipe:
\begin{enumerate}
  \item build $\mC$ and $\mM$ (cotangent weights, mixed or barycentric areas);
  \item choose a total time $t$ of a few $\bar\ell^2$ for noise removal, larger for fairing;
  \item solve $(\mM-t\mC)\bx'=\mM\bx$ once, or a few times with smaller steps;
  \item rescale to the original volume, or use Taubin's filter when only the uniform operator is
  available;
  \item keep features: smoothing blurs sharp edges, which require anisotropic or feature-aware
  variants (for example weights that decrease across large normal changes).
\end{enumerate}

\begin{notebook}{Tema04\_LaplaceBeltrami\_Suavizado.ipynb}
Implements \texttt{uniform\_laplacian}, \texttt{cotan\_matrix} (finite elements),
\texttt{mass\_matrix} and \texttt{mean\_curvature\_normal}; computes the spectral basis with
\texttt{laplacian\_eigs} and a \texttt{spectral\_lowpass} filter; implements
\texttt{smooth\_explicit} with its stability limit and \texttt{smooth\_implicit}; and studies
shrinkage with \texttt{signed\_volume} and \texttt{rescale\_volume}.
\end{notebook}

\begin{keypoints}
\begin{itemize}[leftmargin=*]
  \item The uniform Laplacian depends only on connectivity and drifts tangentially; the cotangent
  Laplacian $\mM^{-1}\mC$ comes from linear finite elements, has linear precision and converges.
  \item Laplacian eigenvectors are a Fourier basis on the surface; smoothing is low-pass filtering.
  \item Explicit diffusion is stable only for $h\le2/\lambda_{\max}$; implicit diffusion,
  $(\mM-h\mC)\bx'=\mM\bx$, is unconditionally stable and approximates $e^{-t\lambda}$.
  \item Diffusion shrinks shapes; volume rescaling and Taubin's $\lambda|\mu$ filter compensate it.
\end{itemize}
\end{keypoints}

\section*{Exercises}
\begin{exercise}
Show that the rows of $\mC$ sum to zero and that $\bfun\T\mC\bfun=-\sum_{ij\in E}C_{ij}(f_i-f_j)^2$.
Conclude that $\mC$ is negative semi-definite when all weights are non-negative.
\end{exercise}
\begin{exercise}
For the uniform Laplacian prove that the eigenvalues of $-\mL_u$ lie in $[0,2]$. Hint: $\mL_u$ is
similar to the symmetric matrix $\mD^{-1/2}\mA\mD^{-1/2}-\mI$.
\end{exercise}
\begin{exercise}
Find the pass-band frequency and the maximum gain of Taubin's filter for $\lambda=0.5$,
$\mu=-0.53$.
\end{exercise}
\begin{exercise}
Under mean curvature flow a sphere satisfies $\dot r=-2/r$. Derive $r(t)=\sqrt{r_0^2-4t}$ and the
collapse time.
\end{exercise}
